\section{ANTECEDENTES}
\label{sec:antecedentes}

La investigación sobre la seguridad de la votación electrónica se ha desarrollado, durante la última década, principalmente en torno a la verificación formal de protocolos sustentados en criptografía de clave pública clásica, como RSA y ElGamal. A partir de 2024, con la publicación de los primeros estándares de criptografía post-cuántica por parte del NIST, este campo comienza a extenderse hacia el análisis de protocolos que incorporan primitivas resistentes a la computación cuántica, un desplazamiento aún incipiente tanto en la literatura internacional como en los antecedentes producidos en el contexto boliviano.

En el ámbito internacional, Baloglu, Bursuc, Mauw y Pang (2021) analizaron la verificabilidad electoral de los sistemas Helios y Belenios mediante verificación formal automatizada con la herramienta Tamarin Prover. El objetivo fue evaluar sistemáticamente propiedades de integridad, confidencialidad y autenticidad bajo distintos escenarios de corrupción. Los resultados permitieron formalizar nuevas garantías de seguridad y detectar ataques no identificados previamente en versiones desplegadas de dichos sistemas, lo que evidencia la utilidad de la verificación formal como método de análisis en protocolos de votación electrónica.

En el ámbito nacional, y particularmente en la Universidad Mayor de San Andrés, Fernández Tristán (2021) propuso un sistema de votación electrónica sustentado en tecnología blockchain, orientado a mejorar la seguridad y la transparencia del registro de votos. En un contexto universitario boliviano distinto al de la UMSA, Churata Sonco (2020) desarrolló, en la Universidad Pública de El Alto, un sistema de votación electrónica basado en protocolos blockchain combinados con cifrado RSA, aplicado como caso de estudio en un proceso electoral universitario. Ambos trabajos comparten el uso de blockchain como mecanismo de integridad del registro de votos, mientras que la protección criptográfica de los datos permanece basada en esquemas clásicos de clave pública, sin considerar la vulnerabilidad de dichos esquemas frente a la computación cuántica.

En el plano de las revisiones de literatura, Ramos y Gómez (2023) realizaron una revisión sistemática de protocolos de votación electrónica basados en criptografía post-cuántica, mediante la cual sistematizaron y compararon las propuestas reportadas en la literatura internacional sobre esta línea de investigación. Dicho trabajo se circunscribe a la síntesis del estado del arte y no propone, diseña ni verifica formalmente un protocolo propio, ni aborda su aplicación en el contexto universitario boliviano; constituye, por tanto, un antecedente de carácter secundario que evidencia el interés creciente por el área, pero no cubre el vacío relativo al diseño y verificación formal de un protocolo específico de votación electrónica en dicho contexto.

El National Institute of Standards and Technology (2024) publicó los estándares FIPS 203 y FIPS 204, referencia técnica que, a la fecha, no ha sido aplicada en los antecedentes nacionales revisados. No se han identificado, dentro de la literatura consultada, antecedentes recientes que aborden específicamente la aplicación de criptografía post-cuántica al diseño de protocolos de votación electrónica, lo que constituye el espacio de investigación en el que se sitúa el presente trabajo.